\chapter{R1 unbounded-$q$: exact contact circuits and their square budgets}
\label{component:p6db37c4233da}
\noindent\textbf{Component }\path{p6db37c4233da}\par
{\small\noindent Source: \path{REC/c144/r1_qo_contact_circuit_budgets.tex}\par}
{\small\noindent SHA256: \path{6db37c4233da1f389bc7a081fa5bad5e3daa4f5d61640483efcd549a455ac952}\par}
{\small\noindent Inventory use: RB1-RB17\par}
\medskip
\section{R1 unbounded-$q$: exact contact circuits and their square budgets}
\label{p6db37c4233da:sec:r1-qo-contact-circuit-budgets}

\paragraph{Status and exact scope.}
This is a collection of proved algebraic lemmas for the remaining
original R1 source sector. It is not a local or all-testing admission.
The complete original parent and the two complete state normals are
RQ6--RQ11 of \path{r1_unbounded_q_structural_checkpoint.tex}.
The new lemmas below avoid substituting $u=vz/T$ and therefore introduce
no inverse powers of a mixed graph coordinate. The final full-normal
allocation is still a separate obligation.

\subsection{The full $n$ scalar has just two negative monomials}

Use RQ1--RQ11. Restore the original $n$ packet in the scalar graph
$N=M=0$ and write this full scalar as $F_{\rm full}$. Equivalently
\[
 F_{\rm full}=F+n(1+v^2)\xi_g^2,\qquad
 \xi_g=p(1+ru^2)v-(1-r)uz.
 \tag{RB1}
\]
Here the factor $1+v^2$ is essential. Its omission would change the
whole scalar and every subsequent allocation.
Put $K_0=CH^2+(D+\tfrac12)(1-H)^2$. Direct expansion gives
\begin{align*}
F_{\rm full}={}&
 Au^2+\{\mu(1-p)+(n+\tfrac12)p^2\}v^2+\alpha z^2\\
&+\{A+\alpha(1-p)^2+\mu(1-p)+\tfrac12(1-p)^2\}u^2v^2
 +nu^2z^2\\
&+\{K_0+AH^2+2H\alpha+H-\tfrac12\}v^2z^2
 +(n+\tfrac12)p^2v^4\\
&+(K_0+n)u^2v^2z^2
 +H^2(A+C+\alpha+\tfrac12)v^4z^2
 +CH^2u^2v^4z^2\\
&-2\{AH+\alpha(1-p)+np+\tfrac12\}uvz\\
&-2\{AH+H\alpha(1-p)+np+
                  \tfrac12[H(1-p)+p]\}uv^3z .
\end{align*}
\par\noindent\hfill\textup{(RB2)}\par
Every displayed even coefficient is nonnegative on the original
source. The only coefficient requiring an estimate is that of
$v^2z^2$: $K_0$ has a positive lower bound of order $1/t$, whereas
the possible negative part is at most $1/2$. In particular this
observation allows arbitrary $n\ge0$, including $n=0$ and unbounded
$n$. It is not a claim that a polynomial with two negative monomials
automatically has a circuit certificate.

\subsection{An exact finite contact expansion}

Let $x=u_c^2$, $y=v_c^2$, and $q_z=z_c^2$ be the actual three
positive contact coordinates of RQ5. Define the separated rows
\[
 U=u/\sqrt{x},\qquad V=v/\sqrt{y},\qquad Z_*=z/\sqrt{q_z}.
 \tag{RB3}
\]
These are rows, not bounded Gram coefficients. Define
\[
 Q_*=U-VZ_*,\qquad B_*=V-UZ_*,\qquad R_*=Z_*-UV,
\]
\[
 \Psi=U^2+V^2+Z_*^2+U^2V^2Z_*^2-4UVZ_* .
 \tag{RB4}
\]
On the positive testing core $\Psi$ is a separated SOS:
\[
 \Psi=(U-V)^2+Z_*^2(1-UV)^2
                     +2UV(1-Z_*)^2.
 \tag{RB5}
\]
The last square has root $\sqrt{2UV}(1-Z_*)$; this is a product
of pure-source and pure-testing roots. Signs and literal testing
axes must be supplied by the inherited core exits in any admission;
the following identities themselves are polynomial identities.

Let $M_a,M_b,M_c,M_{ab},M_{ac},M_{bc},M_{b^2},
M_{abc},M_{b^2c},M_{ab^2c}$ denote the ten positive monomial
coefficients after RB2 is expressed in $(U,V,Z_*)$. Thus, for example,
\[
 M_a=Ax,\quad M_b=\{\mu(1-p)+(n+\tfrac12)p^2\}y,\quad
 M_c=\alpha q_z,\quad M_{ac}=nxq_z,
\]
\[
 M_{b^2c}=H^2(A+C+\alpha+\tfrac12)y^2q_z,\qquad
 M_{ab^2c}=CH^2xy^2q_z.
\tag{RB6}
\]
All coefficients are actual source quantities.
The actual scalar contact value and its first jets imply
\[
 M_b+M_{b^2}=M_{ac}+M_{abc}+M_{ab^2c},\qquad
 M_a+M_{ab}=M_c+M_{bc}+M_{b^2c}.
\tag{RB7}
\]
Set
\[
 k_c=\tfrac12(M_a-M_b+M_{ac}+M_{ab^2c}-M_{bc}),\qquad
 r_c=M_c-M_b+M_{ac}-k_c,
\]
\[
 f_c=M_b-M_{ac},\quad q_c=M_a-f_c,\quad
 e_c=M_{b^2}-M_{ab^2c},\quad d_c=M_{b^2c}-k_c,\quad
 g_c=M_{ab^2c}.
\tag{RB8}
\]
Then the exact full scalar is
\begin{align*}
F_{\rm full}={}&q_cQ_*^2+M_{ac}B_*^2+r_cR_*^2
 +d_cV^2Q_*^2+e_cV^2B_*^2\\
&+f_c\Psi+g_cV^2\Psi
 +k_cZ_*^2(V^2-1)^2 .
\end{align*}
\par\noindent\hfill\textup{(RB9)}\par
The two odd coefficients in RB9 are fixed by the contact value and
the $V$ derivative; RB7 fixes the remaining even coefficients.
No positivity of $e_c,k_c,f_c$ is inferred merely from this identity.
There is also a one-parameter variant, obtained by adding zero:
for any pure-source $\tau$, replace
\[
 (q_c,M_{ac},d_c,e_c,f_c,g_c,k_c)
 \longmapsto
 (q_c-\tau,M_{ac}-\tau,d_c+\tau,e_c+\tau,
  f_c+\tau,g_c-\tau,k_c-\tau)
\]
and add $\tau U^2Z_*^2(V^2-1)^2$.
\par\noindent\hfill\textup{(RB10)}\par

\subsection{Two exact circuit budgets}

The first budget is the polynomial identity
\[
 (1+U^2)\Psi
 =R_*^2+B_*^2+U^2Q_*^2+U^2(1-UVZ_*)^2.
\tag{RB11}
\]
All permutations of $(U,V,Z_*)$ are allowed. In particular
\[
 (1+V^2)\Psi-V^2B_*^2
       =Q_*^2+R_*^2+V^2(1-UVZ_*)^2
\tag{RB12}
\]
is literally a sum of three squares. Consequently a signed
$e_cV^2B_*^2$ term is paid at infinitesimal Gram cost whenever
$e_c^-/f_c$ is source infinitesimal and
$e_c^-/(f_cy)$ is bounded: the cost is
$(e_c^-/f_c)+(e_c^-/(f_cy))v^2$. This criterion never treats
$V$ as a bounded coefficient.

For the second budget put $t_*=UV$, $\delta_*=U-V$, and
\[
 E=Z_*(1-t_*),\qquad T_*=t_*(1-Z_*),\qquad
 J=T_*-\delta_*^2Z_*.
\]
There are exact identities
\[
 R_*=E-T_*,\qquad
 V^2Q_*^2+U^2B_*^2-T_*^2
                =J^2+2t_*\delta_*^2Z_*^2.
\tag{RB13}
\]
Together with RB5 they give the explicit separated SOS remainder
\begin{align*}
&2\{\Psi+V^2Q_*^2+U^2B_*^2\}-R_*^2\\
&\hspace{5mm}=(E+T_*)^2+2\delta_*^2+
 4t_*(1-Z_*)^2+2J^2+4t_*\delta_*^2Z_*^2 .
\end{align*}
\par\noindent\hfill\textup{(RB14)}\par
Thus $R_*^2$ can be bought from the circuit and the two current rows,
without any inverse or root of a mixed sum.
For positive source masses $q_c,v_c,f_c$, a multiplier $a_c\ge0$
has infinitesimal cost if
\[
 a_c/f_c\in\mathrm I,\qquad
 a_c/(q_cy)\text{ is bounded},\qquad
 a_c/(v_cx)\text{ is bounded}.
\tag{RB15}
\]
The last two costs carry the actual testing factors $v^2,u^2$.
The identities continue to hold on literal axes. No division by
$U,V,Z_*$ occurs anywhere in RB11--RB15.

\subsection{The exact remaining $n$ cancellation}

Write $j=\sqrt{xq_z/y}$, so the original actual contact gives
\[
 \frac pj=\frac{1-r}{1+rx},\qquad
 \xi_g=p\sqrt y\,\{B_*-r\sqrt x\,uR_*\}.
\tag{RB16}
\]
Restoring or removing the full $n$ packet must retain the factor
$1+v^2$ from RB1. For its unweighted part alone,
\[
 nxq_z-np^2y=nxq_z\{1-(p/j)^2\},
\]
\[
 n\{1-(p/j)^2\}
 =nr\,\frac{(2-r)+2x+rx^2}{(1+rx)^2}.
\tag{RB17}
\]
This is bounded through the opening $r=0$ in residue and includes
literal $n=0$. The remaining cross and $v^2$ terms are exact
polynomials obtained from RB1 and RB16. They have not been silently
discarded. A proof of their allocation, together with the complete
state normals RQ10--RQ11 and the outer-testing joins, is necessary
before an unbounded-$q$ R1 admission can be recorded.
